% This is samplepaper.tex, a sample chapter demonstrating the
% LLNCS macro package for Springer Computer Science proceedings;
% Version 2.21 of 2022/01/12
%
\documentclass[runningheads]{llncs}
%
\usepackage[T1]{fontenc}
% T1 fonts will be used to generate the final print and online PDFs,
% so please use T1 fonts in your manuscript whenever possible.
% Other font encondings may result in incorrect characters.
%
\usepackage{graphicx}
% Used for displaying a sample figure. If possible, figure files should
% be included in EPS format.
%
% If you use the hyperref package, please uncomment the following two lines
% to display URLs in blue roman font according to Springer's eBook style:
%\usepackage{color}
%\renewcommand\UrlFont{\color{blue}\rmfamily}
%\urlstyle{rm}
%
\usepackage{listings}
\usepackage{xcolor}
\usepackage[utf8]{inputenc}
\usepackage{listings}

% Configuración del estilo para HTML
\lstset{
    language=HTML,
    basicstyle=\ttfamily\small,
    keywordstyle=\color{blue}\bfseries,
    stringstyle=\color{orange},
    commentstyle=\color{gray}\itshape,
    numbers=left,
    numberstyle=\tiny\color{gray},
    stepnumber=1,
    numbersep=8pt,
    backgroundcolor=\color{white},
    showspaces=false,
    showstringspaces=false,
    frame=single,
    rulecolor=\color{black!30},
    tabsize=2,
    breaklines=true,
    breakatwhitespace=true,
    captionpos=b % Leyenda debajo del código
}

\usepackage[backend=biber,style=apa]{biblatex} % puedes cambiar el estilo
\addbibresource{Bibliografia.bib} % enlaza tu archivo .bib

\begin{document}
%
\title{Modulo 1 - Mate Amargo}
%
%\titlerunning{Abbreviated paper title}
% If the paper title is too long for the running head, you can set
% an abbreviated paper title here
%
\author{Emiliano Jaime\inst{1}\orcidID{} \and
Mashael Allay\inst{1}\orcidID{} \and
Estefania Funes Brown\inst{1}\orcidID{}\and
Matias Zaragoza\inst{1}\orcidID{}\and
Nancy Sara\inst{1}\orcidID{}}
%
\authorrunning{M. Zaragoza}
% First names are abbreviated in the running head.
% If there are more than two authors, 'et al.' is used.
%
\institute{Universidad Nacional de Cuyo, Instituto de Ingeniería Inndustrial POBOX 5502 %\and
%Springer Heidelberg, Tiergartenstr. 17, 69121 Heidelberg, Germany
\email{mnzargo28@gmail.com}\\
\url{http://ingenieria.uncuyo.edu.ar} %\and
%ABC Institute, Rupert-Karls-University Heidelberg, Heidelberg, Germany\\
%\email{\{abc,lncs\}@uni-heidelberg.de}}
%
}
\maketitle          %para hacer el titulo    % typeset the header of the contribution
%
\begin{abstract}
Durante el mes de marzo se trabajaron distintos aspectos vinculados al uso de GitHub y la construcción de páginas web académicas. El 4 de marzo se abordó la creación de páginas individuales en GitHub utilizando Markdown, lo que permitió a cada estudiante desarrollar su propio espacio digital. Además, se estableció el vínculo hacia la página grupal, favoreciendo la integración de los trabajos individuales en un entorno colaborativo. Posteriormente, el 11 de marzo se profundizó en el trabajo colaborativo dentro de GitHub, incorporando tablas y la incrustación de imágenes para enriquecer los contenidos. Se introdujo también el uso de HTML bajo los estándares del W3C, lo que brindó herramientas para la construcción de una landing page más profesional y estructurada. Finalmente, el 25 de marzo se destinó a la redacción del informe correspondiente al Módulo 1, empleando la plantilla provista por la cátedra como demo. Esta instancia permitió sistematizar los aprendizajes previos y dar forma a un documento académico coherente y bien organizado. En conjunto, las actividades de marzo consolidaron competencias técnicas en la edición de contenido digital, la colaboración en entornos de desarrollo y la producción de informes académicos, integrando teoría y práctica en un proceso progresivo de aprendizaje.  

\keywords{markup language \and UNCuyo \and TyHM}
\end{abstract}
%
%
%
\section{Uso básico de Github}
GitHub es una plataforma que facilita el control de versiones y la colaboración en proyectos de software. Su uso básico comienza creando un repositorio, donde se almacenan archivos y su historial de cambios. Los desarrolladores pueden clonar el repositorio en sus equipos, realizar modificaciones y subirlas mediante commits y push. Para colaborar, se emplean branches que permiten trabajar en paralelo sin afectar la versión principal. Las propuestas de cambios se gestionan con pull requests, que otros revisan antes de integrarlas. Además, GitHub ofrece herramientas para seguimiento de problemas y documentación, convirtiéndose en un entorno esencial para equipos.  

\subsection{Códigos}
\textbf{Títulos:} se crean con \# al inicio de la línea. 
\newline\textbf{Líneas:} una línea horizontal se hace con --- o *. 
\newline\textbf{Tablas:} se construyen con barras verticales | y guiones -. \newline Ejemplo:  
  \newline| Columna 1 | Columna 2 |
  \newline|\texttt{----------------}|\texttt{----------------}|
  \newline| Dato A    | Dato B    |
\newline\textbf{Listas:} 
  \begin{itemize}
      \item Ordenadas: 1. Primer ítem
      \item No ordenadas: - Ítem o * Ítem
  \end{itemize}
\textbf{Links:} [texto](URL) crea un enlace.  
\newline\textbf{Hipervínculos:} igual que links, se pueden insertar dentro del texto: Visita [GitHub](https://github.com).  
\newline\textbf{Letras (estilos):}  
  \begin{itemize}
      \item Negrita: *texto*
      \item Cursiva: texto
      \item Tachado: \texttt{}texto\texttt{}
  \end{itemize}
GitHub usa Markdown para dar formato a títulos, listas, tablas, enlaces y estilos de texto en archivos como README.md.

\subsection{Uso básico de Markdown}

Markdown es un lenguaje de marcado ligero que facilita dar formato a texto. Se usa para crear títulos con \#,
listas con - o , y enfatizar palabras con *cursiva o *negrita*. Los enlaces se escriben como [texto](URL) 
y las imágenes como ![alt](URL). Los bloques de código se delimitan con ```, mientras que las citas se marcan 
con >. También permite tablas con | y separadores con ---. Su ventaja es la simplicidad: el texto sigue siendo 
legible sin procesar y se convierte fácilmente en HTML, ideal para documentación, notas y contenido web.

\subsection{Citas}
\subsubsection{Citas de Github}
\begin{itemize}
    \item \cite{lopez2003latex}
    \item \cite{kalliamvakou2014github}
    \item \cite{dabbish2012social}
    \item \cite{githubdocs2025}
    \item \cite{ray2014github}
    \item \cite{tsay2014github}
\end{itemize}
\subsubsection{Citas de Markdown}
\begin{itemize}
    \item \cite{gruber2004markdown}
    \item \cite{macfarlane2014pandoc}
    \item \cite{perkel2017markdown}
    \item \cite{stackexchange2020markdown}
    \item \cite{markdownguide2025}
\end{itemize}

%\subsubsection{Sample Heading (Third Level)} Only two levels of headings should be numbered. Lower level headings remain unnumbered; they are formatted as run-in headings.

%\paragraph{Sample Heading (Fourth Level)}
%The contribution should contain no more than four levels of headings. Table~\ref{tab1} gives a summary of all heading levels.
\section{Uso básico de HTML}
HTML (HyperText Markup Language) es el lenguaje estándar para estructurar páginas web. Su uso básico consiste en organizar el contenido mediante etiquetas que delimitan elementos. Por ejemplo, <h1> define títulos, <p> párrafos y <a> enlaces. Las listas se crean con <ul> y <ol>, mientras que las tablas se construyen con <table>, <tr> y <td>. Las imágenes se insertan con <img> y los hipervínculos con <a href="URL">. HTML permite jerarquizar la información y combinar texto, multimedia y vínculos, sirviendo como la base sobre la cual se aplican estilos con CSS y funcionalidades con JavaScript.  

\subsection{Códigos}
\begin{lstlisting}[caption={Estructura básica de un documento HTML5.}, label={list:html_ejemplo}]
<html>
<head>
  <h1>Mi primer pagina web</h1>
</head>
<body>
  <p>Hola, mi gente estamos en el internet</p>
  <ul>
    <li>Aguante River</li>
    <li>Colapinto</li>
    <li>y Messi</li>
  </ul>
</body>
</html>
\end{lstlisting}

\subsection{Citas}
\begin{itemize}
    \item \cite{w3c1991html}
    \item \cite{musciano2012html}
    \item \cite{mdnhtmlcite}
    \item \cite{w3schoolshtmlcite}
    \item \cite{w3cscholarlyhtml}
\end{itemize}

\section{Uso básico de LaTex}
LaTeX es un sistema de composición tipográfica muy usado en documentos científicos y técnicos. 
Su uso básico consiste en escribir texto plano con comandos que controlan formato y estructura. 
Se inicia con \verb|\documentclass{}| para definir el tipo de documento, seguido de 
\verb|\begin{document}| y \verb|\end{document}|. 
Los comandos como \verb|\section{}|, \verb|\textbf{}| o \verb|\emph{}| permiten organizar y dar estilo. 
Para matemáticas, se emplean entornos como \verb|$...$| para expresiones en línea o \verb|\[...\]| para ecuaciones destacadas. 
LaTeX separa contenido de diseño, generando documentos profesionales y consistentes con mínima intervención manual.

\subsection{Códigos}
\textbf{Estructura:} todo documento comienza con \texttt{\textbackslash documentclass\{...\}} y termina con \texttt{\textbackslash end\{document\}}.  
\newline\textbf{Secciones:} se crean con \texttt{\textbackslash section\{Título\}}, \texttt{\textbackslash subsection\{...\}}.  
\newline\textbf{Texto:} se puede dar formato con \texttt{\textbackslash textbf\{negrita\}}, \texttt{\textbackslash textit\{cursiva\}}.  
\newline\textbf{Matemáticas:} se escriben entre \texttt{\$ ... \$} para inline o con entornos como \texttt{\textbackslash begin\{equation\}}.   
\newline\textbf{Listas:}  
\begin{itemize}
    \item Ordenadas: \texttt{\begin{enumerate} \item Primer ítem \end{enumerate}}
    \item No ordenadas: \texttt{\begin{itemize} \item Ítem \end{itemize}}
\end{itemize}
\textbf{Tablas:} se construyen con \texttt{\\begin{tabular}\{...\}} usando \texttt{|} y \texttt{\&}.  

{\begin{itemize} \item Ejemplo \end{itemize}}
\begin{lstlisting}[caption={Estructura básica de un documento LaTeX.}, label={list:latex_ejemplo}]
\documentclass{article}
\usepackage[utf8]{inputenc}

\begin{document}

\section{Mi primer documento}
Este es un ejemplo sencillo escrito en LaTeX.

\subsection{Texto destacado}
Podemos usar \textbf{negrita} o \textit{cursiva} facilmente.

\subsection{Matematicas}
Una formula: \( E = mc^2 \).

\end{document}
\end{lstlisting}


\subsection{Citas}
\begin{itemize}
    \item \cite{lamport1994latex}
    \item \cite{mittelbach2004latexcompanion}
    \item \cite{overleaflatex}
    \item \cite{manualdelatex2020}
    \item \cite{aprendeconalf2020}
\end{itemize}

\section{Uso básico de Google Colab}
Google Colab es una plataforma en la nube que permite trabajar con notebooks interactivos más allá de Python. 
Aunque su entorno principal está optimizado para este lenguaje, también admite R, Julia y otros mediante kernels personalizados o extensiones. 
Se pueden ejecutar comandos de sistema con \verb|!| y scripts en C/C++ o JavaScript con configuraciones adicionales. 
Los notebooks combinan texto, código y resultados, y se integran con Google Drive para guardar y colaborar en tiempo real. 
Además, Colab ofrece acceso gratuito a GPU y TPU, instalación de paquetes con \verb|pip| o \verb|apt-get|, y gran flexibilidad para proyectos de programación y ciencia de datos.

\subsection{Códigos}
\textbf{Celdas de código:} se escriben en Python y se ejecutan con \texttt{Shift+Enter}.  
\newline\textbf{Celdas de texto:} usan Markdown para títulos, listas, enlaces y estilos.  
\newline\textbf{Importación de librerías:} se hace con \texttt{import}, por ejemplo: \texttt{import numpy as np}.  
\newline\textbf{Gráficos:} se generan con librerías como \texttt{matplotlib}.  
\newline\textbf{Ejemplo básico:}  
\newline\texttt{import numpy as np}  
\newline\texttt{import matplotlib.pyplot as plt}  
\newline\texttt{x = np.linspace(0, 10, 100)}  
\newline\texttt{y = np.sin(x)}  
\newline\texttt{plt.plot(x, y)}  
\newline\texttt{plt.show()}  
\newline\textbf{Listas:}  
\begin{itemize}
    \item Ordenadas: 1. Primer ítem
    \item No ordenadas: - Ítem o * Ítem
\end{itemize}
\textbf{Hipervínculos:} se insertan en celdas de texto con Markdown, por ejemplo:  
Visita [Google Colab](https://colab.research.google.com).  


\subsection{Citas}
\begin{itemize}
    \item \cite{bisong2019colab}
    \item \cite{googlecolabdocs}
    \item \cite{googleblog2017colab}
    \item \cite{carneiro2018colab}
    \item \cite{stackoverflow2020colab}
\end{itemize}

%
% the environments 'definition', 'lemma', 'proposition', 'corollary',
% 'remark', and 'example' are defined in the LLNCS documentclass as well.
%
%
% ---- Bibliography ----
%
% BibTeX users should specify bibliography style 'splncs04'.
% References will then be sorted and formatted in the correct style.
%
% \bibliographystyle{splncs04}
% \bibliography{mybibliography}
%

\end{document}               
